\begin{tutorial}{Impact on Polaris Deep}

Process the impacts in reverse order: instead of deleting junctions, add them back one by one.

Start from the state after all $q$ impacts and count the active junctions $V$, the connections $E$ between them, and their connected components $C$ directly. Keep the components in a disjoint-set-union structure.

When a junction is restored, add its vertex ($V$ and $C$ increase by $1$), then look at its at most four neighbours. For every active neighbour such that the restored junction and the neighbour open towards each other, add the corresponding edge ($E$ increases by $1$); if the two endpoints were in different components, unite them ($C$ decreases by $1$).

After each restoration the number of independent loops equals $E - V + C$. Recording these values while time runs backwards and printing them in reverse order answers every impact. Each edge is considered at most once, so the running time is $O((n \cdot m + q)\, \alpha(n \cdot m))$ per test case.

Forward simulation does not work directly: deleting one junction can change the number of loops by an amount that depends on how its neighbours are connected after its removal, which is expensive to determine online.

(Why does counting with the formula $E - V + C$ match the definition in the legend? For a spanning forest of the network, every non-tree connection closes exactly one loop containing a connection no previous loop used, giving $E - V + C$ admissible loops; conversely, any loop with a fresh connection can be charged to that connection, so no sequence can be longer.)


\end{tutorial}
